% u01-polyglossia: polyglossia with three languages and scripts (Latin,
% Greek, Cyrillic) in TeX Live's fonts, hyphenation per language
\documentclass{article}
\usepackage{fontspec}
\usepackage{polyglossia}
\setdefaultlanguage{english}
\setotherlanguages{french,greek,russian}
\setmainfont{Libertinus Serif}
\begin{document}
\section{Languages}
This paragraph is in English, with its own hyphenation patterns and the
usual spacing after a full stop. It goes on long enough for the line breaker
to hyphenate a word or two in a narrow measure of text.

\begin{french}
Ce paragraphe est en français : la ponctuation haute prend une espace
fine insécable, et les guillemets « comme ceci » aussi. La césure suit les
règles du français dans une mesure assez étroite.
\end{french}

\begin{greek}
Αυτή η παράγραφος είναι στα ελληνικά, με τους δικούς της κανόνες
συλλαβισμού και τη σωστή στίξη· το κείμενο συνεχίζεται αρκετά ώστε να
χωριστεί σε γραμμές.
\end{greek}

\begin{russian}
Этот абзац написан по-русски, с русскими правилами переноса и
типографики. Текст продолжается достаточно долго, чтобы разбиться на
несколько строк.
\end{russian}

Back in English: \textfrench{une citation}, \textgreek{μια φράση} and
\textrussian{фраза} inline.
\end{document}
